% PDF pp. 17--22; continuation of proof of 4.4.
This reduces us to the case where $A$ is local and noetherian. Let us now distinguish two cases.

1$^\circ$) $A$ is local noetherian and henselian. Let $\widehat A$ be its completion, $\widehat S=\Spec(\widehat A)$ and $\widehat H=H\times_S\widehat S$. Apply theorem 3.7: one finds a $\widehat S$-group $\widehat G$ of multiplicative type, isotrivial and of finite type, and a homomorphism of $\widehat S$-groups
\[
\widehat u:\widehat G\to\widehat H
\]
which is an open immersion and a closed immersion, such that $\widehat u$ induces an isomorphism $\widehat u_0:\widehat G_0\isomto\widehat H_0$.

By remark 4.0.1, the base change functor by $\widehat S\to S$ induces an equivalence between the category of isotrivial groups of multiplicative type over $S$ and over $\widehat S$; in particular $\widehat G$ ``comes from'' an $S$-group of multiplicative type $G$, isotrivial and of finite type. By 4.3, $\widehat u$ comes from a homomorphism
\[
u:G\to H;
\]
moreover, $u$ is an open and closed immersion and induces an isomorphism $u_0:G_0\to H_0$, since this is so after the faithfully flat quasi-compact base change $\widehat S\to S$. This therefore proves 4.4 in this case (taking of course, in the conclusion of 4.4, $S'=S$ and $s'=s$).

2$^\circ$) $A$ is local noetherian. Reduction to case 1$^\circ$) is immediate by applying 1$^\circ$) to the ``henselization'' $A^h$ of $A$. More precisely, one sees easily (using SGA 1, I \S5\footnote{Or EGA IV$_4$, \S18.6.}) that the local rings $\OO_{S',s'}$ of étale $S$-preschemes $S'$ endowed with a point $s'$ over $s$ such that $\kappa(s')/\kappa(s)$ is trivial form a filtering inductive system whose inductive limit is a noetherian henselian local ring $A^h$ (called the ``henselization'' of $A$); for details of this construction (due to Nagata in the normal case), cf. SGA 4, VIII \S4.\footnotemark[\value{footnote}] The sorites of EGA IV$_3$, \S8 then allow one, as in part a) of the proof, to deduce from a result known over the inductive limit $A^h$ of the $\OO_{S',s'}$ an analogous result over one of the $(S',s')$, which proves 4.4.
\begin{corollary}{4.5}\label{X.4.5}
Let $S$ be a prescheme, $H$ an $S$-group prescheme of multiplicative type and of finite type. Then $H$ is quasi-isotrivial, i.e. is split by a surjective étale morphism $S'\to S$.
\end{corollary}
Indeed, let $s\in S$. By 4.4, there exists an étale morphism $S'\to S$, an $s'\in S'$ over $s$ such that $\kappa(s)=\kappa(s')$, and a subgroup $G'$ of $H'$, of isotrivial multiplicative type and of finite type, such that $G'_{s'}=H'_{s'}$. Since $G',H'$ are of multiplicative type and of finite type, by IX 2.9 there exists an open neighborhood $U'$ of $s'$ such that $G'|U'=H'|U'$.

\nde{The original has been expanded in the preceding and following passages.} Suppose in addition $S$ local henselian with closed point $s$; then by EGA IV$_4$, 18.5.11 b) there exists a section $\sigma$ of $S'\to S$ such that $\sigma(s)=s'$. (N.B. One may see directly that $B=\OO_{S',s'}$ equals $A=\OO_{S,s}$ as follows: by 4.3.0 (i), one has $B\simeq A/I$, and since $B$ is an $A$-algebra of finite presentation and flat, $I$ is an ideal of finite type (cf. EGA IV$_1$, 1.4.7), and $I=I\goth m$ (where $\goth m$ is the maximal ideal of $A$), whence $I=0$.) Thus $H$ is already isotrivial. One therefore obtains:
\begin{corollary}{4.6}\label{X.4.6}
Let $A$ be a henselian local ring, $k$ its residue field, and $\pi$ the topological Galois group of an algebraic closure of $k$.

(i) Every group of multiplicative type and of finite type over $S=\Spec(A)$ is isotrivial.

(ii) The category of these groups over $S$ is equivalent (via the functor $H\mto H\otimes_A k$) to the analogous category over $S_0=\Spec(k)$; it is therefore anti-equivalent, by 1.4, to the category of Galois modules under $\pi$ of finite type as $\ZZ$-modules.
\end{corollary}
\begin{corollary}{4.7}\label{X.4.7}
Under the conditions of 4.4, suppose in addition one of the following conditions satisfied:

a) For every generization $t$ of $s$, $H_t$ is of multiplicative type and of the same type as $H_s$.

b) $H$ is commutative and the ${}_nH$ ($n>0$) are finite over $S$ in a neighborhood of $s$.

c) For every generization $t$ of $s$, the fiber $H_t$ is connected.

Then there exists an open neighborhood $U$ of $s$ such that $H|U$ is of multiplicative type.
\end{corollary}
Taking account of the fact that an étale morphism is open, one is reduced to proving that there exists (with the notation of the conclusion of 4.4) an open neighborhood $U'$ of $s'$ such that $G'|U'=H'|U'$. Put $S''=\Spec(\OO_{S',s'})$; as $G',H'$ are of finite presentation over $S'$, it suffices to show, by EGA IV$_3$, 8.8.2, that $G''=H''$. One may therefore suppose $S=S''$, and then hypotheses (a), (b), (c) become those of the lemma below, whose proof is the same as that of 3.8:
\begin{lemma}{4.7.1}\label{X.4.7.1}
Let $S$ be a local scheme, $s$ its closed point, $H$ an $S$-group scheme of finite type, $G$ an open and closed subgroup of $H$, of multiplicative type, such that $G_s=H_s$. Suppose in addition one of the following conditions satisfied:

a) The fibers of $H$ are of multiplicative type and of the same type as $H_s$.

b) $H$ is commutative and the ${}_nH$ ($n>0$) are finite over $S$.

c) The fibers of $H$ are connected.

Then $G=H$.
\end{lemma}
\begin{corollary}{4.8}\label{X.4.8}
Let $S$ be a prescheme, $H$ an $S$-group prescheme affine, flat and of finite presentation over $S$, with fibers of multiplicative type.

For $H$ to be of multiplicative type, it is necessary and sufficient that it satisfy one of the following two conditions (equivalent under the preceding conditions):

a) The type of $H_s$ (cf. IX 1.4) is a locally constant function of $s\in S$.

b) $H$ is commutative, and the ${}_nH$ ($n>0$) are finite over $S$.

Moreover, these conditions a), b) are implied by the following:

c) The fibers of $H$ are connected.
\end{corollary}
In particular, one finds:
\begin{corollary}{4.9}\label{X.4.9}
Let $S$ be a prescheme, $H$ an $S$-group prescheme flat and of finite presentation over $S$. Suppose in addition $H$ affine over $S$ with connected fibers.\nde{This result is generalized in 8.1: in fact it suffices to suppose the fibers of $H$ affine and connected.} If $s\in S$ is such that $H_s$ is a torus, there exists an open neighborhood $U$ of $s$ such that $H|U$ is a torus.

In particular, if all the fibers of $H$ are tori, $H$ is a torus.
\end{corollary}

\sgasection{5}{Scheme of homomorphisms from one group of multiplicative type into another. Twisted constant groups and groups of multiplicative type}
\begin{definition}{5.1}\label{X.5.1}
a) Let $S$ be a prescheme, $R$ a group prescheme over $S$. One says that $R$ is a \emph{twisted constant group over $S$} if it is locally isomorphic, in the sense of the faithfully flat quasi-compact topology, to a constant group scheme, i.e. of the form $M_S$ with $M$ an ordinary group.

b) One says that the twisted constant group $R$ over $S$ is \emph{quasi-isotrivial}, resp. \emph{isotrivial}, resp. \emph{locally isotrivial}, resp. \emph{locally trivial}, resp. \emph{trivial}, if in the above definition one can replace the faithfully flat quasi-compact topology by the étale topology, resp. the global finite étale topology, resp. the finite étale topology, resp. the Zariski topology, resp. the coarsest (or ``chaotic'') topology, cf. IX 1.1 and 1.2, and IV 6.6.

Saying that $R$ is quasi-isotrivial (resp. isotrivial) therefore means that there exists a surjective étale (resp. and finite) morphism $S'\to S$ such that $R'=R\times_S S'$ is a constant group over $S$; saying that it is trivial means that $R$ is a constant group.
% typo?: "constant group over S" rather than S' retained as printed.
c) One defines as in VIII 1.4 the \emph{type} of a twisted constant group $R$ over $S$ at an $s\in S$; it is an isomorphism class of ordinary groups which, for variable $s$, is a locally constant function of $s$, and hence constant if $S$ is connected. One will also say that $R$ is ``of type $M$'' if all the fibers of $R$ are of type $M$.

One should take care that $R$ is quasi-compact over $S$ only if it is finite over $S$, i.e. if its type at every $s\in S$ is a finite group.\nde{See lemma 5.12 below.}

d) The case of greatest interest to us is that where $R$ is commutative. We shall then say that $R$ is ``finitely generated'' if its type at every point $s\in S$ is given by a $\ZZ$-module of finite type, a notion which must not be confused with the schematic notion ``$R$ of finite type over $S$'' (cf. above).
\end{definition}
\begin{remark}{5.2}\label{X.5.2}
We shall also have to consider $S$-preschemes $X$ which are locally isomorphic (for the faithfully flat quasi-compact topology) to constant preschemes, independently of any group structure. We shall then say that $X$ is a \emph{twisted constant bundle over $S$}, and extend to these preschemes the terminology introduced in 5.1. Of course, one should take care that when $X$ is endowed with an $S$-group structure, the meaning of the expressions ``twisted constant'', ``isotrivial'', etc. changes according to whether one takes account of the group structure over $S$. The same is true if one considers on $X$ any other kind of algebraic structure (for example that of a Galois principal bundle, which will be considered in the next \No).
% typo: "platequasi-compacte" -> "flat quasi-compact".
\end{remark}
\begin{proposition}{5.3}\label{X.5.3}
Let $R$ be a commutative twisted constant group over $S$.

(i) The functor $H=D_S(R)$ (cf. VIII 1) is representable and is a group of multiplicative type over $S$.

(ii) For every $s\in S$, the type of $R$ at $s$ equals that of $H$ at $s$.

(iii) For $R$ to be quasi-isotrivial (resp. isotrivial, resp. trivial, resp. locally isotrivial, resp. locally trivial), it is necessary and sufficient that $H$ be so.
\end{proposition}
\begin{remark}{5.3.1}\label{X.5.3.1}
\nde{This remark, which appeared after the proof in the original, has been moved here.} One may be concerned to see the term ``type'' used in 5.3 in two different senses according to whether $R$ or $H$ is involved; fortunately, when an $S$-group prescheme $G$ is both a twisted constant group and of multiplicative type, its type in either sense is the same, thanks to the fact that an ordinary finite commutative group is isomorphic to its dual!
\end{remark}
\begin{proof}
As covering families for the faithfully flat quasi-compact topology are families of effective descent for the fibered category of affine group preschemes over variable base preschemes (SGA 1, VIII 2.1), one sees that $H$ is representable (and affine over $S$), since it is so ``locally''\nde{Cf. VIII \S1.7.} (for it is so when $R$ is constant, and then $H$ is a diagonalizable group).

The fact that $H$ is of multiplicative type is then evident by definition, as is the fact that the type of $R$ and $H$ at $s\in S$ is the same. Finally, since $H_{S'}=D_{S'}(R_{S'})$, the last assertion is reduced to the ``trivial'' case, i.e. to checking that $R$ is trivial if and only if $H$ is so, which follows immediately from biduality theorem VIII 1.2.
\end{proof}
To finish specifying the correspondence between twisted constant groups and groups of multiplicative type, one must start with a group of multiplicative type $H$ and study $R=D_S(H)$. If the latter is representable, it is evidently a twisted constant group, and one shall have $H\simeq D_S(R)$. In other words:
\begin{scholium}{5.4.0}\label{X.5.4.0}
The functor $R\mto D_S(R)$ is an anti-equivalence\nde{``Equivalence'' has been corrected to ``anti-equivalence''.} between the category of twisted constant groups over $S$ and that of groups of multiplicative type $H$ over $S$ such that $D_S(H)$ is representable.
\end{scholium}
I do not know whether this condition on $H$ is always satisfied; we shall nevertheless see that it is when $H$ is quasi-isotrivial, in particular when $H$ is of finite type.
\begin{lemma}{5.4}\label{X.5.4}
Let $S'\to S$ be a faithfully flat morphism locally of finite presentation, $X'$ an $S'$-prescheme separated, locally of finite presentation and locally quasi-finite over $S'$.

Then every descent datum on $X'$ relative to $S'\to S$ is effective, i.e. there exists an $S$-prescheme $X$ and an $S'$-isomorphism $X\times_S S'\simeq X'$ compatible with the descent datum.
\end{lemma}
One has already noted that the hypothesis on $S'\to S$ implies that it is a covering morphism for the faithfully flat quasi-compact topology (IV 6.3.1 (iv)), a fortiori a morphism of effective descent.

When $X'\to S'$ is quasi-compact, hence of finite presentation and quasi-finite, it is a quasi-affine morphism (cf. SGA 1, VIII 6.2\nde{(When $S'$ is locally noetherian, and EGA IV$_3$, 8.11.2 in general.)}), and effectiveness follows in this case from SGA 1, VIII 7.9. In the general case one reduces immediately to the case where $S,S'$ are affine. One covers $X'$ by affine open subsets $U'_i$; let $V'_i$ be the saturation of $U'_i$ for the equivalence relation in $X'$ defined by the descent datum, i.e. $V'_i=q_2(q_1^{-1}(U'_i))$, where $q_1,q_2$ are the two projections of $X''_1=X'\times_{S'}(S'',p_1)$ onto $X'$ ($q_1=\operatorname{pr}_1$, and $q_2$ is deduced from the first projection of $X''_2=X'\times_{S'}(S'',p_2)$ by the given descent isomorphism $X''_1\simeq X''_2$). Since $S'\to S$ is faithfully flat quasi-compact locally of finite presentation, so is $p_1:S''=S'\times_S S'\to S'$, hence so are $q_1,q_2$, which are consequently open morphisms (SGA 1 IV 6.6\nde{(When $S'$ is locally noetherian, and EGA IV$_2$, 2.4.6 in general.)}). Consequently $V'_i$ is an open and quasi-compact subset of $X'$. By what has already been seen, the descent data induced on the $V'_i$ are effective, whence it follows that the datum on $X'$ is also so (SGA 1, VIII 7.2).
\begin{corollary}{5.5}\label{X.5.5}
A faithfully flat morphism of finite presentation $S'\to S$ is a morphism of effective descent for the fibered category of twisted constant groups (over variable base preschemes).
\end{corollary}
Indeed, this amounts to asserting the effectiveness of a descent datum under the conditions of 5.4 when $X'$ is a constant $S'$-prescheme.
\begin{theorem}{5.6}\label{X.5.6}
Let $S$ be a prescheme, $G,H$ two $S$-group preschemes of quasi-isotrivial multiplicative type, with $G$\nde{$H$ has been corrected to $G$.} of finite type.

Then $\uHom_{S\text{-}\mathrm{gr.}}(H,G)$ is representable by, and is, a quasi-isotrivial twisted constant group over $S$;\nde{The following has been added.} for every $s\in S$, if the type of $G$ (resp. $H$) at $s$ is $M$ (resp. $N$), that of $\uHom_{S\text{-}\mathrm{gr.}}(H,G)$ is $\Hom_{\mathrm{gr.}}(M,N)$.
\end{theorem}
One proceeds as in 4.2, using the fact that the assertion is established (VIII 1.5) when $G,H$ are trivial. The necessary effectiveness criterion is supplied by 5.5 (in the case of a surjective étale morphism $S'\to S$).

In particular, putting $G=\mathbf{G}_{\mathrm{m},S}$, one finds:
\begin{corollary}{5.7}\label{X.5.7}
(i) Let $H$ be an $S$-group of quasi-isotrivial multiplicative type; then the $S$-group $D_S(H)$ is representable and is a quasi-isotrivial twisted constant group over $S$.

(ii) The functors $R\mto D_S(R)$ and $H\mto D_S(H)$ are anti-equivalences, quasi-inverse to one another, between the category of quasi-isotrivial twisted constant $S$-groups and that of $S$-groups of quasi-isotrivial multiplicative type.

(iii) These functors induce anti-equivalences between the subcategories consisting of isotrivial, resp. locally isotrivial, resp. locally trivial, resp. trivial groups.
\end{corollary}
The last assertion is mentioned only for the record, being already contained in 5.3.

Moreover, since every $S$-group of multiplicative type and of finite type is quasi-isotrivial by 4.5, one deduces from 5.6:
\begin{corollary}{5.8}\label{X.5.8}
Let $S$ be a prescheme, $G,H$ two $S$-group preschemes of multiplicative type and of finite type; then $\uHom_{S\text{-}\mathrm{gr.}}(H,G)$ is representable and is a finitely generated quasi-isotrivial twisted constant $S$-group over $S$.
\end{corollary}
Note also that in 5.3, $R$ is finitely generated if and only if $H=D_S(R)$ is of finite type (IX 2.1 b)). By 4.5, $H$ is then quasi-isotrivial, so $R$ is quasi-isotrivial. One thus finds:
\begin{corollary}{5.9}\label{X.5.9}
The functors of 5.7 induce anti-equivalences quasi-inverse to one another between the category of $S$-groups $H$ of multiplicative type of finite type and that of finitely generated twisted constant $S$-groups $R$; moreover, every such group $R$ is quasi-isotrivial.
\end{corollary}
\begin{corollary}{5.10}\label{X.5.10}
Let $H,G$ be two $S$-group preschemes of multiplicative type and of finite type.

\begingroup\emergencystretch=3em
(i) Then $\underline{\Isom}_{S\text{-}\mathrm{gr.}}(H,G)$ is representable by an open and closed subprescheme of $\uHom_{S\text{-}\mathrm{gr.}}(H,G)$, and is a twisted constant $S$-prescheme.\par
\endgroup

(ii) In particular, $\underline{\Aut}_{S\text{-}\mathrm{gr.}}(H)$ is representable and is a twisted constant $S$-group (in general noncommutative).
\end{corollary}
This follows from 5.8 and VIII 1.6.\nde{Correct this reference by treating the case of the functor $\underline{\Isom}_{S\text{-}\mathrm{gr.}}(H,G)$ in an addition 1.5.1 in VIII\dots}
\begin{enoncerm}{Recollection}{5.11.0}\label{X.5.11.0}
\nde{This recollection has been added; it is used several times in what follows. (See also EGA I, 6.1.9; note, however, that the proof in loc. cit. seems unnecessarily complicated.)} Recall that if $X$ is a locally noetherian prescheme, its connected components (which are always closed) are open. Indeed, let $C$ be a connected component of $X$, $x\in C$ and $U$ a noetherian open neighborhood of $x$; then $U$ has only finitely many connected components, so the connected component $U_x$ of $x$ in $U$ is open in $U$, hence in $X$; as $U_x\subseteq C$, this shows that $C$ is open in $X$.
\end{enoncerm}
\begin{proposition}{5.11}\label{X.5.11}
Let $S$ be a prescheme, $R$ a commutative twisted constant group over $S$, $H=D_S(R)$ the group of multiplicative type it defines. Consider the following conditions:

(i) $H$ is isotrivial (i.e. $R$ is isotrivial).

(ii) $R$ is a union of subpreschemes $R_i$ which are both open and closed and quasi-compact over $S$ (and then necessarily finite over $S$).
% Proposition continues on PDF p. 23, in en-05.tex.
